% =====================================================================
%  CUERPO DEL MANUAL DEL ADMINISTRADOR - PharmaWeigh
%  Capturas reales del sistema. Sin credenciales.
% =====================================================================

\seccion{1. Acceso al sistema}

\tablaAcceso

\vspace{8pt}

\reglasAcceso

\captura{01_login.png}{Pantalla de entrada. No expone usuarios ni contraseñas de ejemplo.}

\captura{19_login_error.png}{Credenciales incorrectas. El mensaje es siempre el mismo, no revela si falló el correo o la contraseña.}

\begin{cajaOjo}
\textbf{Tu cuenta abre todo el sistema.} Da de alta usuarios, libera lotes y firma
fases. Trátala como una cuenta privilegiada: contraseña larga, no compartida, y un
usuario nominal para cada persona (nada de una cuenta \textit{admin} de equipo).
\end{cajaOjo}

\avisoValidacion

% =====================================================================
\newpage
\seccion{2. Alcance del rol}

Como \manualRol{} ves \textbf{las ocho} opciones del menú y tienes todos los permisos
que el sistema concede. Lo importante es saber qué \textbf{no existe} para nadie:

\begin{center}
\begin{tabularx}{\linewidth}{@{}>{\raggedright\arraybackslash}X c@{}}
\toprule
\textbf{Acción} & \textbf{Existe} \\
\midrule
Crear usuarios, órdenes, recetas, materiales y lotes  & Sí \\
Aprobar, rechazar y retener lotes                     & Sí \\
Registrar pesajes y firmar fases                      & Sí \\
Consultar la bitácora                                 & Sí \\
\midrule
\textbf{Cambiar o reiniciar una contraseña}           & \textbf{No, desde la pantalla} \\
\textbf{Dar de baja o desactivar un usuario}          & \textbf{No, desde la pantalla} \\
\textbf{Editar o borrar un usuario}                   & \textbf{No} \\
\textbf{Editar o borrar una receta}                   & \textbf{No} \\
\textbf{Cancelar una orden}                           & \textbf{No} \\
\textbf{Corregir o borrar un pesaje}                  & \textbf{No} \\
\textbf{Editar o borrar un lote}                      & \textbf{No} \\
\textbf{Borrar o editar la bitácora}                  & \textbf{No, ni desde el servidor} \\
\textbf{Exportar la bitácora}                         & \textbf{No, desde la pantalla} \\
\bottomrule
\end{tabularx}
\end{center}

\begin{cajaNota}
Las tres primeras filas de la segunda mitad (contraseñas, bajas y exportación) se
resuelven \textbf{desde la consola del servidor}. El procedimiento está en el documento
de operación del sistema, \cd{docs/OPERACION.md}. La sección 4 de este manual explica
lo esencial.
\end{cajaNota}

\captura{03_dashboard.png}{Dashboard: conteos del día, procesos activos, últimas órdenes y actividad reciente.}

% =====================================================================
\newpage
\seccion{3. Usuarios y roles}

\subseccion{3.1 Crear un usuario}

\begin{enumerate}
    \item \menu{Usuarios} $\rightarrow$ \boton{+ Nuevo Usuario}.
    \item \campo{Nombre}: nombre completo real. Es el que va a aparecer en cada renglón
          de la bitácora y en cada firma.
    \item \campo{Email}: es el identificador de entrada. Único; si se repite, el sistema
          responde \cd{Ya existe un registro con esos datos.}
    \item \campo{Contraseña}: \textbf{mínimo 10 caracteres}. Es la \textbf{única} regla
          que el sistema impone: no exige mayúsculas, números ni símbolos, no caduca y no
          guarda historial.
    \item \campo{Rol}: ver la tabla de la sección 3.3.
    \item \campo{Badge (código de barras)}: opcional. Si lo pones, el gafete aparece en
          \menu{Códigos de Barras} para imprimirlo.
    \item \boton{Crear Usuario}.
\end{enumerate}

\captura{21_usuario_nuevo.png}{Formulario de alta. Bajo el campo de contraseña se indica el mínimo de 10 caracteres.}

\vspace{6pt}

Los errores se muestran en pantalla, sobre el formulario, y \textbf{no se crea nada}:

\captura{28_error_formulario.png}{Correo ya registrado: el sistema lo rechaza y conserva lo que habías capturado.}

\begin{cajaOjo}
\textbf{Entrega la contraseña por un canal aparte} y pide que se cambie. Pero ojo:
\textbf{el sistema no tiene pantalla de cambio de contraseña}. En la práctica, cambiarla
significa que tú la regeneres desde el servidor. Tenlo en cuenta al definir el
procedimiento de la planta.
\end{cajaOjo}

\newpage
\subseccion{3.2 La lista de usuarios}

\captura{18_usuarios.png}{Nombre, correo, rol, gafete y estado de cada cuenta.}

\vspace{6pt}

La tabla es informativa: \textbf{no hay botones de editar, borrar ni desactivar}. La
columna \menu{ESTADO} muestra \textit{Activo} o \textit{Inactivo}, pero se cambia desde
el servidor, no desde aquí. Supervisión ve esta misma lista, sin el botón de alta.

\subseccion{3.3 Qué ve y qué puede cada rol}

\begin{center}
\footnotesize
\begin{tabularx}{\linewidth}{@{}l >{\raggedright\arraybackslash}X@{}}
\toprule
\textbf{Rol} & \textbf{Menú y capacidades} \\
\midrule
\textbf{ADMIN}      & Las 8 opciones. Todo lo que el sistema permite, incluido crear usuarios y firmar. \\
\textbf{SUPERVISOR} & Las 8 opciones. Todo menos crear usuarios (los ve, no los crea). Firma. \\
\textbf{OPERARIO}   & Dashboard, Órdenes, Dispensado, Recetas, Inventario, Códigos. Pesa. No firma, no aprueba, no crea nada. \\
\textbf{CALIDAD}    & Dashboard, Órdenes, Dispensado, Recetas, Inventario, Auditoría. Aprueba/rechaza/retiene lotes y \textbf{firma}. No pesa, no crea nada. \\
\textbf{DESARROLLO} & Dashboard, Órdenes, Dispensado, Recetas, Inventario, Auditoría. \textbf{Crea recetas}. No pesa, no firma, no aprueba. \\
\textbf{ALMACEN}    & Sólo Dashboard, Inventario y Códigos. Crea materiales, recibe lotes y \textbf{retiene}; \textbf{no aprueba ni rechaza}. \\
\textbf{AUDITOR}    & Dashboard, Órdenes, Dispensado, Recetas, Inventario, Auditoría. Sólo lectura, en todo. \\
\bottomrule
\end{tabularx}
\end{center}

\vspace{6pt}

Los permisos salen de una \textbf{matriz única} que aplican por igual el menú, las
pantallas y el servidor. \textbf{Falla cerrada}: lo que no está concedido explícitamente,
está negado. Cambiar el menú de un rol no es una opción de configuración: es un cambio
de código.

\begin{cajaTip}[Separación de funciones]
Los roles ya vienen separados a propósito: quien recibe no libera (Almacén no aprueba),
quien pesa no firma (Operario no firma), quien diseña no produce (Desarrollo no crea
órdenes). Si le das \textbf{ADMIN} a media planta, esa separación desaparece. Asigna el
rol más restringido que permita hacer el trabajo.
\end{cajaTip}

% =====================================================================
\newpage
\seccion{4. Lo que se hace desde el servidor}

Cuatro operaciones habituales \textbf{no tienen pantalla} y se ejecutan en consola, con
acceso al servidor. El procedimiento completo, con los comandos exactos, está en
\cd{docs/OPERACION.md}.

\begin{center}
\begin{tabularx}{\linewidth}{@{}p{0.34\linewidth} >{\raggedright\arraybackslash}X@{}}
\toprule
\textbf{Necesidad} & \textbf{Cómo se resuelve} \\
\midrule
Reponer una contraseña olvidada & Desde consola. No hay ``olvidé mi contraseña'' ni correo de recuperación en la aplicación. \\
Dar de baja a alguien que sale de la empresa & Desde consola, marcando la cuenta como inactiva. Una cuenta inactiva no puede entrar. \\
Crear el primer administrador de una instalación & Con el guion \cd{crear-usuario}, que también sirve de rescate si se pierde el acceso. \\
Extraer la bitácora para una auditoría o un respaldo & Consulta directa a la base de datos. \\
\bottomrule
\end{tabularx}
\end{center}

\begin{cajaOjo}
\textbf{Un usuario bloqueado por intentos fallidos no se desbloquea a mano.} Son 5
fallos con el mismo correo en 15 minutos y el bloqueo dura 15 minutos. No hay comando
ni botón para levantarlo antes: se espera. Dilo así cuando te llamen de piso.
\end{cajaOjo}

\begin{cajaNota}
Las altas hechas por consola quedan en la bitácora como \cd{CREAR\_USUARIO\_CONSOLA},
distinguibles de las hechas por pantalla (\cd{CREAR\_USUARIO}). Es la forma de auditar
las intervenciones administrativas.
\end{cajaNota}

% =====================================================================
\newpage
\seccion{5. Inventario}

\subseccion{5.1 El ciclo de vida de un lote}

\begin{itemize}
    \item \textbf{Todo lote entra en \cuarentena}, sin excepción.
    \item \textbf{Aprobar} y \textbf{rechazar}: Calidad, Supervisión y Administración.
    \item \textbf{Retener} (regresar un lote aprobado a cuarentena): los anteriores
          y además Almacén.
    \item \rechazado{} y \agotado{} son \textbf{finales}. \agotado{} lo pone el sistema
          al llegar a cero.
    \item \textbf{No se puede recibir un lote ya caducado.}
\end{itemize}

\begin{cajaNota}
\textbf{Criterio de caducidad, por confirmar con el cliente.} La fecha es de calendario
de la planta (sin hora ni desfase por huso horario) y el lote \textbf{deja de poder
usarse a las 00:00 de su día de caducidad}. Si la planta usa otro criterio, hay que
resolverlo antes del arranque: cambia qué lotes acepta la estación el día del corte.
\end{cajaNota}

\captura{15_inventario_lotes.png}{Tabla de lotes con la columna de acciones.}

\begin{cajaNota}
\boton{Aprobar} y \boton{Retener} se aplican con \textbf{un solo clic} (son reversibles
entre sí). \boton{Rechazar} pide \textbf{confirmación en dos pasos} dentro de la propia
fila --- \cd{¿Rechazar LOT-...?}, nombrando el lote, con \boton{Sí, rechazar} y
\boton{Cancelar}, y escondiendo el \boton{Aprobar} de esa fila --- porque \rechazado{}
es un estado final. Conviene mencionarlo en la capacitación de Calidad y Supervisión:
la confirmación protege del clic accidental y del clic en la fila equivocada, pero
\textbf{no de la decisión}: una vez confirmada no hay deshacer.
\end{cajaNota}

\subseccion{5.2 Materiales y recepción}

\boton{+ Nuevo Material} pide código de barras, nombre, unidad, stock mínimo y
descripción. \textbf{La unidad no se puede cambiar después} y el material no se puede
borrar: es el dato que conviene revisar dos veces en un alta.

\captura{16_inventario_recepcion.png}{Recepción: se identifica el material por su código y después se capturan los datos del lote.}

% =====================================================================
\newpage
\seccion{6. Recetas y órdenes}

\subseccion{6.1 Recetas}

Las cargan Desarrollo, Supervisión o Administración. Estructura: receta
$\rightarrow$ fases $\rightarrow$ ingredientes. El orden es el orden real de pesaje y
\textbf{cada fase exige una firma} que bloquea la siguiente.

\begin{cajaOjo}
\textbf{Ninguna receta se edita ni se borra}, y el campo de versión no lo incrementa
ninguna pantalla. Un cambio de fórmula significa dar de alta otra receta con código
distinto. Documenta esto en el procedimiento de la planta antes del arranque: es la
limitación que más sorprende a los usuarios.
\end{cajaOjo}

\subseccion{6.2 Órdenes}

\captura{20_orden_nueva.png}{Alta de orden: receta, lote de producto final, multiplicador y prioridad.}

\vspace{6pt}

El \campo{multiplicador} escala el target \textbf{y el rango de tolerancia} de cada
ingrediente. Los estados avanzan solos: \pendiente{} $\rightarrow$ \enproceso{} con el
primer pesaje, y $\rightarrow$ \dispensado{} con la firma de la última fase.

\textbf{No hay cancelación de órdenes desde la pantalla.} Una orden creada por error se
queda \pendiente{} para siempre, ensuciando la lista de dispensado, salvo que se
intervenga la base de datos.

% =====================================================================
\newpage
\seccion{7. Dispensado y firma}

Puedes entrar a la estación y también firmar. Las reglas que impone el servidor, y que
vas a tener que explicar cuando llamen de piso:

\begin{itemize}
    \item \textbf{La báscula no está conectada.} El operario \textbf{teclea} el peso.
          El sistema valida el rango, no el instrumento ni su calibración.
    \item \textbf{Un peso fuera de tolerancia no se registra.} Se rechaza con mensaje.
          No existe una autorización de excepción en ninguna pantalla.
    \item \textbf{Los pasos van en orden estricto}, y una fase completa bloquea la
          siguiente hasta que se firme.
    \item \textbf{El doble clic no duplica} el pesaje ni el descuento de inventario.
    \item \textbf{Nunca hay stock negativo.}
    \item \textbf{Sólo lotes aprobados y no caducados} se pueden dispensar.
    \item \textbf{Un pesaje confirmado no se corrige ni se borra} desde ninguna pantalla.
\end{itemize}

\captura{22_firma_fase.png}{Firma de fase: el firmante escribe su propio correo y su propia contraseña.}

\subseccion{Reglas de la firma}

\begin{itemize}
    \item Firman \textbf{SUPERVISOR}, \textbf{CALIDAD} y \textbf{ADMIN}, y nadie más.
    \item El rechazo responde siempre el mismo mensaje, sin decir qué falló.
    \item Cada rechazo queda en bitácora como \cd{FIRMA\_RECHAZADA}.
    \item \textbf{5 rechazos} con el mismo correo en 15 minutos bloquean ese correo
          15 minutos, también para entrar al sistema.
\end{itemize}

% =====================================================================
\newpage
\seccion{8. Bitácora}

\menu{Auditoría} muestra los \textbf{500 eventos más recientes}, del más nuevo al más
viejo, con fecha y hora, usuario y rol, acción, entidad, \textbf{IP} y detalles.

\captura{17_auditoria.png}{Bitácora con filtros por acción y por usuario.}

\vspace{6pt}

Eventos registrados:

\vspace{4pt}
\noindent\begin{tabularx}{\linewidth}{@{}XXX@{}}
\cd{LOGIN}            & \cd{CREAR\_USUARIO}       & \cd{CREAR\_ORDEN} \\
\cd{LOGOUT}           & \cd{CREAR\_MATERIAL}      & \cd{DISPENSAR} \\
\cd{LOGIN\_FALLIDO}   & \cd{RECEPCION\_LOTE}      & \cd{FIRMAR\_FASE} \\
\cd{FIRMA\_RECHAZADA} & \cd{CAMBIAR\_ESTADO\_LOTE}& \cd{FIRMAR\_DISPENSADO} \\
                      & \cd{CREAR\_RECETA}        & \cd{COMPLETAR\_DISPENSADO} \\
\end{tabularx}

\begin{cajaInfo}[Inmutabilidad]
La bitácora \textbf{no se puede editar ni borrar}, y la restricción vive en la base de
datos, no en la aplicación. Tu cuenta tampoco puede. Si alguna vez hiciera falta
depurarla, sería una intervención deliberada sobre la base, no una operación normal
del sistema.
\end{cajaInfo}

\begin{cajaOjo}
\textbf{Sin filtro por fecha, sin búsqueda y sin exportación.} Cuando Auditoría o
Calidad pidan evidencia de un periodo, la extracción la haces tú desde el servidor.
Deja ese procedimiento acordado antes de la primera inspección, no durante.
\end{cajaOjo}

% =====================================================================
\newpage
\seccion{9. Códigos de barras}

\menu{Códigos de Barras} genera las etiquetas en pantalla, en tres pestañas:
\menu{Lotes}, \menu{Materiales} y \menu{Gafetes}. La de gafetes sólo la ven
Administración y Supervisión, y lista únicamente a los usuarios activos que tengan un
gafete asignado.

\captura{26_codigos.png}{Página de códigos, lista para imprimir o para abrir en un celular.}

\vspace{6pt}

No hay botón de imprimir dentro de la aplicación: se usa la impresión del navegador
(Ctrl+P o Cmd+P); la página está preparada para salir legible en blanco y negro.
La otra forma de usarla es abrirla en un teléfono y escanear desde ahí con la cámara.

% =====================================================================
\newpage
\seccion{10. Equipo y navegador}

\begin{itemize}
    \item \textbf{Navegador:} Chrome o Edge actualizados. La aplicación es web; no se
          instala nada en la estación.
    \item \textbf{Pistola lectora:} cualquiera que funcione como teclado (HID). No
          requiere configuración ni controladores. Lee el código y manda el Enter sola.
    \item \textbf{Cámara:} el campo de escaneo tiene un botón para leer con la cámara del
          equipo o del celular, útil como respaldo si falla la pistola.
    \item \textbf{Báscula:} \textbf{no se integra}. No hay conexión serial, USB ni de red
          con ningún instrumento. El peso se teclea.
    \item \textbf{Pantalla:} las capturas de este manual están tomadas a 1440$\times$900.
          Con menos ancho, la tabla de lotes y la de auditoría requieren desplazamiento
          horizontal.
\end{itemize}

% =====================================================================
\newpage
\seccion{11. Qué sí puedes romper}

\begin{center}
\begin{tabularx}{\linewidth}{@{}p{0.32\linewidth} >{\raggedright\arraybackslash}X@{}}
\toprule
\textbf{Riesgo} & \textbf{Por qué importa} \\
\midrule
\textbf{Repartir el rol ADMIN}      & Destruye la separación de funciones: un mismo usuario podría recibir, liberar, pesar y firmar. Es lo primero que revisa una auditoría. \\
\textbf{Cuentas compartidas}        & La bitácora deja de identificar personas. Toda la trazabilidad del sistema se apoya en que cada cuenta sea de alguien. \\
\textbf{Confirmar un rechazo a la ligera} & La confirmación en dos pasos evita el clic accidental, pero \boton{Sí, rechazar} sigue siendo irreversible. \\
\textbf{Guardar una receta mal}     & No se edita ni se borra; se va a producción tal cual. \\
\textbf{Crear órdenes de prueba en producción} & No se pueden cancelar: se quedan en la lista de dispensado indefinidamente. \\
\textbf{No documentar el procedimiento de contraseñas y bajas} & Son operaciones de consola. Si no están escritas ni asignadas, en la práctica no ocurren. \\
\textbf{Prometer cumplimiento normativo} & El sistema \textbf{no está validado}. Afirmarlo en un documento comprometería a la planta ante una inspección. \\
\bottomrule
\end{tabularx}
\end{center}

% =====================================================================
\newpage
\seccion{12. Preguntas frecuentes}

\subseccion{Un usuario olvidó su contraseña}
No hay recuperación en la aplicación. Tienes que generar una nueva desde el servidor
(\cd{docs/OPERACION.md}) y entregársela por un canal aparte.

\subseccion{Se fue una persona de la empresa}
Su cuenta se desactiva desde el servidor. No hay botón de baja en la pantalla, y una
cuenta que sigue activa sigue pudiendo entrar y firmar.

\subseccion{Alguien se bloqueó por intentos fallidos}
Son 15 minutos y no hay forma de levantarlo antes. Que espere.

\subseccion{Creé una orden por error}
No se puede cancelar desde la pantalla. Se queda \pendiente. La única salida es
intervenir la base de datos.

\subseccion{¿Puedo borrar un pesaje mal registrado?}
No, por diseño. Documenta la desviación por el procedimiento de la planta; el registro
original queda en la bitácora.

\subseccion{Me piden la bitácora de los últimos seis meses}
La pantalla sólo muestra 500 registros y no exporta. La extracción se hace desde el
servidor.

\subseccion{¿El sistema cumple 21 CFR Part 11 o NOM-059?}
No lo afirmes. PharmaWeigh incluye funciones orientadas a la trazabilidad del
dispensado --- bitácora inmutable, firma electrónica, control de tolerancias y de
estados de lote --- pero \textbf{no está validado} y no hay documentación de
calificación que respalde un enunciado de cumplimiento. La validación es un proyecto
aparte, a cargo de la planta.
